% AUTORIDAD DE FUENTE:
% expose_VIII_body.tex, líneas 2367--2446; índices de escaneo 317--321;
% folios de fuente 306--310. La familia P_{eta Q_eta}/P_eta x_eta Q_eta
% se conserva reversiblemente mediante \sisi. La fórmula q_eta^*(E_eta)
% sigue directamente el escaneo y repara solo la transcripción.

\medskip
\noindent 7.3.4. \underline{Reducción de las obstrucciones de tipo
(7.1.2) a las de tipo (7.1.1).}

Supongamos, para fijar ideas, que existe un entero $n>0$ tal que
$n\Phi=0$, y sea
\[
 q\in Q(S)
\]
una sección de $Q$. Por abuso de notación, escribamos
\[
 q_0\in\Psi(k)
\]
para la imagen en $\Psi(k)$ del valor de $q$ en $s$. Considerando
$E_\eta$ como una extensión de
$\sisi{P_{\eta Q_\eta}}{P_\eta\times_\eta Q_\eta}$ por
$\underline G_{mQ_\eta}$ sobre $Q_\eta$, la imagen inversa da una
extensión
\[
 q_\eta^*(E_\eta),\qquad
 \text{extensión de $P_\eta$ por $\underline G_{m\eta}$},
\]
y, por tanto, una obstrucción
\begin{equation}\tag{7.3.4.1}
 d_q\overset{\mathrm{def}}{=}
 d\bigl(q_\eta^*(E_{\sisi{*}{\eta}})\bigr):
 \Phi\longrightarrow(\mathbb Q/\mathbb Z)_k.
\end{equation}
Si $d=d(E_\eta)$, se tiene la relación
\begin{equation}\tag{7.3.4.2}
 d_q(x)=d(x,q_0),
\end{equation}
válida para todo punto $x$ de $\Phi$ con valores en un esquema $T$
sobre $k$.

\smallskip
\noindent 7.3.4.3. Para todo $q\in Q(S)$, esta fórmula determina $d_q$
a partir del emparejamiento $d$; recíprocamente, permite recuperar
manifiestamente $d$ a partir de los invariantes $d_q$, al menos cuando
todo punto de $\Psi$ procede de una sección de $Q$. Por el lema de
Hensel, esta condición se cumple automáticamente cuando $S$ es
estrictamente local. El emparejamiento $d$ queda determinado una vez
conocido sobre la clausura separable de $k$, y su formación conmuta con
la localización étale sobre $S$ por 7.3.5.3. Así, sustituyendo $S$ por
su localización estricta si es necesario, la determinación de $d$ se
reduce a la de los invariantes $d_q$.

En particular, cuando todo punto de $\Psi$ procede de una sección $q$
de $Q$, el emparejamiento $d$ es nulo —es decir, $E_\eta$ se prolonga
a una biextensión $E$ de $(P,Q)$ por $\underline G_m$— si y solo si
$d_q$ es nulo para toda sección $q$ de $Q$, es decir, si y solo si
$q_\eta^*(E_\eta)$ se prolonga a una extensión de $P$ por
$\underline G_m$. Sin hipótesis sobre $\Psi$, el criterio análogo
sigue siendo válido siempre que se permita que los $q$ sean puntos con
valores en esquemas $S'$ étale y locales sobre $S$.

%%%%%%%%%%  scan idx 319  |  folio 308  %%%%%%%%%%

\smallskip
\noindent 7.3.4.4. Recíprocamente, partiendo de la situación de 7.1(a),
la obstrucción (7.1.1) puede interpretarse como una obstrucción de tipo
(7.1.2): tómese $Q=\mathbb Z_S$ y la biextensión de
$(P_\eta,\mathbb Z_\eta)$ deducida de $E_\eta$. En efecto, basta aplicar
la construcción precedente con $q$ igual a la sección $1$ de
$\mathbb Z_S$.

\medskip
\noindent 7.3.5. \underline{Efecto de un cambio de base $S'\to S$.}
Supongamos dado un morfismo local dominante de traits
\[
 f:S'\longrightarrow S,
\]
y denotemos por $\eta'$ el punto genérico de $S'$, por $s'$ su punto
cerrado, y pongamos $k'=k(s')$. Si $S=\operatorname{Spec}(V)$ y
$S'=\operatorname{Spec}(V')$, entonces $V'$ es un anillo de valoración
discreta que domina al anillo de valoración discreta $V$. Pongamos
\begin{equation}\tag{7.3.5.1}
 e(S'/S)=e(V'/V)=
 \operatorname{longitud}_{V'}(V'/\mathfrak m V'),
\end{equation}
donde $\mathfrak m$ es el ideal maximal de $V$.

Sea ahora $P$ como en 7.1(a), sea $E_\eta$ una extensión de $P_\eta$
por $\underline G_{m\eta}$, pongamos $P'=P\times_S S'$, y sea
$E'_{\eta'}$ la extensión de $P'_{\eta'}$ por
$\underline G_{m\eta'}$ obtenida por imagen inversa de $E_\eta$. Con
la notación de 4.1(a), existe un isomorfismo canónico
\[
 \Phi'\overset{\mathrm{def}}{=}P'_0/{P'_0}^{\circ}
 \simeq\Phi\otimes_k k'\simeq\Phi_{k'},
\]
y, por tanto, podemos considerar
\[
 d(E_\eta)_{k'}:\Phi'\longrightarrow(\mathbb Q/\mathbb Z)_{k'}.
\]
Entonces
\begin{equation}\tag{7.3.5.2}
 d(E'_{\eta'})=e(S'/S)\,d(E_\eta)_{k'}.
\end{equation}

%%%%%%%%%%  scan idx 320  |  folio 309  %%%%%%%%%%

Análogamente, bajo las hipótesis de 7.1(b) y con la notación evidente,
(7.3.5.2) sigue siendo válida, denotando ahora $d(E_\eta)$ el
emparejamiento (7.1.2), $d(E_\eta)_{k'}$ el emparejamiento
\[
 \Phi'\otimes\Psi'\longrightarrow(\mathbb Q/\mathbb Z)_{k'}
\]
obtenido de él por el cambio de base $k\to k'$, y $E'_{\eta'}$ la
biextensión de $(P'_{\eta'},Q'_{\eta'})$ por
$\underline G_{m\eta'}$ obtenida de la biextensión $E_\eta$ de
$(P_\eta,Q_\eta)$ por $\underline G_{m\eta}$ mediante el cambio de
base $\eta'\to\eta$.

\smallskip
\noindent 7.3.5.3. Un caso interesante es $e(S'/S)=1$, es decir,
\[
 \mathfrak m'=\mathfrak m V',
\]
donde $\mathfrak m'$ denota el ideal maximal de $V'$. Esto ocurre, en
particular, si $V'$ es étale sobre $V$, o si $V'$ es el henselizado, el
henselizado estricto o el completado de $V$. En este caso, (7.3.5.2)
dice simplemente que \underline{la formación de la obstrucción (7.1.1),
respectivamente (7.1.2), conmuta con el cambio de base considerado
$S'\to S$}.

\smallskip
\noindent 7.3.5.4. Por otra parte, (7.3.5.2) muestra en el caso general
que, para toda extensión $E_\eta$ como en 7.1(a) (respectivamente, toda
biextensión $E_\eta$ como en 7.1(b)), existe un morfismo
\underline{finito} de traits $S'\to S$ tal que $k(\eta')$ sea una
extensión separable de $k(\eta)$ y tal que la extensión
(respectivamente, biextensión) $E'_{\eta'}$ obtenida de $E_\eta$ por
cambio de base $\eta'\to\eta$ se prolongue a una extensión
(respectivamente, biextensión) $E'$ sobre todo $S'$. En efecto, basta
elegir $S'/S$ de modo que $e(S'/S)$ sea múltiplo del entero $n>0$ que
aparece en 7.1, y siempre puede encontrarse tal extensión con
$k(\eta')$ separable sobre $k(\eta)$.

%%%%%%%%%%  scan idx 321  |  folio 310  %%%%%%%%%%
